\documentclass[11pt,a4paper]{article}

\usepackage[margin=1in]{geometry}
\usepackage{listings}
\usepackage{xcolor}
\usepackage{tikz}
\usetikzlibrary{shapes.geometric,arrows.meta,positioning,fit,backgrounds}
\usepackage{booktabs}
\usepackage{hyperref}
\usepackage{caption}
\usepackage{amsmath}
\usepackage{enumitem}
\usepackage{fancyhdr}
\usepackage{titling}
\usepackage{float}

\hypersetup{colorlinks=true, linkcolor=blue, urlcolor=blue, citecolor=blue}

\definecolor{codebg}{rgb}{0.96,0.96,0.96}
\definecolor{codekw}{rgb}{0.0,0.0,0.6}
\definecolor{codecomment}{rgb}{0.3,0.5,0.3}
\definecolor{codestring}{rgb}{0.6,0.0,0.0}

\lstdefinestyle{base}{
  backgroundcolor=\color{codebg},
  basicstyle=\ttfamily\small,
  commentstyle=\color{codecomment}\itshape,
  keywordstyle=\color{codekw}\bfseries,
  stringstyle=\color{codestring},
  numbers=left,
  numberstyle=\tiny\color{gray},
  numbersep=8pt,
  frame=single,
  rulecolor=\color{gray!40},
  breaklines=true,
  breakatwhitespace=false,
  showstringspaces=false,
  tabsize=2,
  columns=fullflexible,
  keepspaces=true
}
\lstset{style=base}

\lstdefinelanguage{pseudo}{
  morekeywords={if,else,for,in,repeat,return,raise,function,not,every,delete,goto,while,True,False,continue},
  morecomment=[l]{//},
  sensitive=true
}
\lstdefinestyle{pseudostyle}{
  language=pseudo,
  backgroundcolor=\color{codebg},
  basicstyle=\ttfamily\small,
  commentstyle=\color{codecomment}\itshape,
  keywordstyle=\color{codekw}\bfseries,
  frame=single,
  rulecolor=\color{gray!40},
  breaklines=true,
  breakatwhitespace=false,
  showstringspaces=false,
  tabsize=2,
  columns=fullflexible,
  keepspaces=true
}

\pagestyle{fancy}
\fancyhf{}
\fancyhead[L]{Function Calling in the APARA C Compiler}
\fancyhead[R]{\thepage}
\renewcommand{\headrulewidth}{0.4pt}

\setlength{\droptitle}{-1.5cm}
\title{\textbf{Function Calling in the APARA C Compiler}\\[2mm]
\large How C Function Calls Are Lowered to APARA mcode}
\author{APARA C Compiler Project}
\date{\today}

\begin{document}
\maketitle
\thispagestyle{fancy}
\vspace{-8mm}

% ============================================================================
\section{Introduction and Strategy}
% ============================================================================

A C function call touches more compiler machinery at once than almost any other
construct: the IR generator (one node, \texttt{IRCall}), the register allocator (every
live value must survive a call whose callee may clobber anything), and the code
generator's calling-convention logic (arguments in, return value out, the stack frame
itself). This report covers all three, end to end, for direct calls; indirect calls
through a function pointer are mentioned only briefly, since they are blocked at the
assembler level for an unrelated reason (Section~\ref{sec:limits}).

\textbf{The governing strategy} has three deliberate simplifications, each trading some
performance for a much smaller, more predictable amount of code to get right:
\begin{enumerate}
  \item \textbf{Fixed argument registers, no stack-passed arguments.} Exactly four
        registers, \texttt{\$r2}--\texttt{\$r5}, carry the first four arguments; a 5th
        argument is a hard compile-time error, not a silent truncation or a fallback to
        stack passing.
  \item \textbf{Full caller-save, not a fixed callee-save/caller-save split.} Every
        register the allocator considers \emph{currently live} at a call site --- not a
        fixed subset chosen by an ABI table --- is spilled to the stack before the call
        and reloaded after. The callee is therefore free to clobber every register
        except \texttt{\$r26}/\texttt{\$r27} (FP/SP, never touched by the allocator at
        all).
  \item \textbf{One epilogue per function, no matter how many \texttt{return}
        statements.} Every \texttt{return} just loads the result into \texttt{\$r1} and
        jumps to a single, shared epilogue label; the actual stack teardown and
        \texttt{\$return} instruction exist exactly once per function.
\end{enumerate}

% ============================================================================
\section{The Call Site (Caller Side): Algorithm}
\label{sec:caller}
% ============================================================================

\texttt{ir\_gen.py}'s \texttt{\_call} lowers an ordinary call to exactly one IR node,
\texttt{IRCall(dest, func\_name, args)} (ir\_gen.py:1188) --- all of the complexity below
is \texttt{codegen.py}'s \texttt{\_gen\_IRCall} (codegen.py:943--1018) turning that one
node into real mcode.

\begin{lstlisting}[style=pseudostyle, caption={\texttt{\_gen\_IRCall}, codegen.py:943--1018 (essence)}]
function EMIT-CALL(dest, func_name, args):
    if len(args) > 4:  COMPILE-ERROR("too many arguments"); exit

    // 1. Caller-save: spill every currently-live register to the stack
    saved <- current allocator map (name -> register), as a list
    for (name, reg) in saved:
        slot <- next caller-save stack slot
        emit "$st [FP + slot], reg"
        remember slot for name

    // 2. Set up arguments in $r2..$r5, reading from the JUST-SAVED stack
    //    slots (not from registers) so writing arg[i] can never clobber
    //    a register that arg[i+1] still needs to read
    for i, arg in enumerate(args[:4]):
        if arg is a live temp that was just saved:  emit "$ld ARG[i], [FP + slot]"
        elif arg is a constant:                      load it into ARG[i] directly
        else:                                          emit "+ ARG[i] = 0 + arg's register"

    // 3. Make the call
    emit "$call func_name"
    // From here on every register except $r1 may hold garbage --
    // the allocator's bookkeeping is now stale; only the caller-save
    // slots on the stack are trustworthy.

    // 4. Capture the return value
    if dest is needed:
        if no free register exists:
            EVICT-ONE-CALLER-SAVED-VALUE-TO-A-SPILL-SLOT()   // see Section 2.1
        reg <- allocate a register for dest
        emit "+ reg = 0 + $r1"

    // 5. Restore every caller-saved value from its stack slot
    for (name, reg, slot) in saved:
        emit "$ld reg, [FP + slot]"
\end{lstlisting}

\subsection*{The one tricky edge case: capturing the return value when the pool is full}
\label{sec:caller-edge}
If all 28 registers were already live before the call, step 4 has nowhere to put the
return value --- every register is still recorded as holding some caller-saved value
(even though, physically, the call just clobbered it). The fix: pick any one saved value
that is \emph{not} currently sitting in \texttt{\$r1}, reload its correct value from its
caller-save slot, write it straight back out to a dedicated long-term spill slot instead,
and mark its register free. This never touches \texttt{\$r1} itself, so the live return
value is never at risk of being overwritten while making room for it.

\begin{figure}[H]
\centering
\begin{tikzpicture}[
  scale=0.82, every node/.style={scale=0.82},
  node distance=6mm and 10mm,
  proc/.style={rectangle, draw, thick, rounded corners, minimum width=5.6cm, minimum height=0.8cm, align=center, fill=orange!10},
  decision/.style={diamond, draw, thick, aspect=2.4, minimum width=4.2cm, minimum height=0.7cm, align=center, fill=blue!8, font=\small},
  term/.style={rectangle, draw, thick, rounded corners, minimum width=4.2cm, minimum height=0.7cm, align=center, fill=green!10},
  arr/.style={-{Triangle[length=2.2mm,width=1.8mm]}, thick}
]
  \node[term] (start) {\texttt{IRCall(dest, func, args)}};
  \node[proc, below=of start] (save) {1. Save every live register to the stack};
  \node[proc, below=of save] (args) {2. Load \$r2..\$r5 from saved slots / constants};
  \node[proc, below=of args] (call) {3. emit \texttt{\$call func}};
  \node[decision, below=of call] (destq) {dest needed?};
  \node[decision, right=22mm of destq] (fullq) {pool full?};
  \node[proc, above=of fullq] (evict) {evict one non-\$r1 saved value to a spill slot};
  \node[proc, below=of fullq] (capture) {allocate reg; copy from \$r1};
  \node[proc, below=of destq] (restore) {5. Reload every saved value from its slot};
  \node[term, below=of restore] (done) {done};

  \draw[arr] (start) -- (save);
  \draw[arr] (save) -- (args);
  \draw[arr] (args) -- (call);
  \draw[arr] (call) -- (destq);
  \draw[arr] (destq) -- node[above,font=\scriptsize]{yes} (fullq);
  \draw[arr] (destq) -- node[left,font=\scriptsize]{no} (restore);
  \draw[arr] (fullq) -- node[right,font=\scriptsize]{yes} (evict);
  \draw[arr] (evict) -- (fullq);
  \draw[arr] (fullq) -- node[right,font=\scriptsize]{no} (capture);
  \draw[arr] (capture) |- (restore);
  \draw[arr] (restore) -- (done);
\end{tikzpicture}
\caption{Call-site code generation, step by step. The eviction loop (top right) only
ever runs once per call, since one evicted register is always enough to make room for the
single return-value destination.}
\end{figure}

% ============================================================================
\section{Function Entry and Exit (Callee Side): Algorithm}
\label{sec:callee}
% ============================================================================

\begin{lstlisting}[style=pseudostyle, caption={\texttt{\_gen\_IRFuncBegin} / \texttt{\_gen\_IRFuncEnd}, codegen.py:572--619 (essence)}]
function EMIT-PROLOGUE(func_name, params, frame_size):
    if len(params) > 4:  COMPILE-ERROR("too many parameters"); exit
    reset register allocator to a fresh 28-register pool   // new function, clean slate

    LABEL func_name
    emit "$st [SP + 0], FP"             // save caller's frame pointer
    emit "+ FP = 0 + SP"                // establish the new frame
    total <- frame_size + CALLER_SAVE_BYTES + SPILL_RESERVE
    emit "- SP = SP - total"            // carve out locals + caller-save + spill areas

    for i, (name, fp_offset, width) in enumerate(params):
        emit "$st (width) [FP + fp_offset], ARG[i]"   // width-correct, not a blind $i64

function EMIT-EPILOGUE(func_name):
    LABEL func_name + "_epilogue"       // every "return" jumps here, never falls through
    emit "+ SP = 0 + FP"                // collapse the frame
    emit "$ld FP, [SP + 0]"             // restore caller's frame pointer
    emit "$return"
\end{lstlisting}

The parameter-store width matters: every parameter is written back at its \emph{real} C
width (\texttt{\$i32} for a plain \texttt{int}, \texttt{\$i8} for a \texttt{char}, \ldots),
not a blanket 64-bit store. APARA's narrow load always reads a fixed bit-slice of the
8-byte DMEM word; storing a small value as a full 64-bit quantity and then reading it back
narrow would silently read the wrong half. (This was a real, confirmed bug in this
compiler before it was fixed --- see Section~\ref{sec:limits}.)

\begin{lstlisting}[style=pseudostyle, caption={\texttt{\_gen\_IRReturn}, codegen.py:1104--1111}]
function EMIT-RETURN(value):
    if value exists:
        emit "+ $r1 = 0 + value's register"     // load the result into RET
    emit unconditional jump to func_name + "_epilogue"
\end{lstlisting}
A function with three \texttt{return} statements compiles to three small
load-into-\texttt{\$r1}-then-jump sequences, all converging on the \emph{same} epilogue
--- the stack-teardown code is never duplicated.

% ============================================================================
\section{A Fully Compiled, Verified Worked Example}
\label{sec:worked}
% ============================================================================

\begin{lstlisting}[language=C, caption={Source}]
int add(int a, int b) { return a + b; }
int main() {
    int r;
    r = add(3, 4);
    return r;
}
\end{lstlisting}

\begin{lstlisting}[caption={Generated mcode (both functions, in full)}]
add:
||
    $st ($i64) [$r27 + 0] $r26        // prologue: save caller's FP
;
||
    + $r26 ($i64) $r0 $r27            // FP = SP
    $set $r1 0 824                    // 824 = frame(88) + caller-save(224) + spill(512)
;
||
    - $r27 ($i64) $r27 $r1            // SP -= 824
    $st ($i32) [$r26 + -8]  $r2       // store param a (width = $i32, matches "int")
    $st ($i32) [$r26 + -16] $r3       // store param b
    + $r1 ($i64) $r26 -8              // body starts: &a
;
||
    $ld ($i32) $r2 [$r1 + 0]          // load a
    + $r3 ($i64) $r26 -16
;
||
    $ld ($i32) $r4 [$r3 + 0]          // load b
;
||
    + $r5 ($i64) $r2 $r4              // r5 = a + b
;
||
    + $r1 ($i64) $r0 $r5              // RET = r5
    ? ($i64) $r0 == $goto add_epilogue
;
add_epilogue:
||
    + $r27 ($i64) $r0 $r26            // SP = FP
;
||
    $ld ($i64) $r26 [$r27 + 0]        // FP = caller's FP
    $return
;
main:
||
    $st ($i64) [$r27 + 0] $r26
;
||
    + $r26 ($i64) $r0 $r27
    $set $r1 0 816
;
||
    - $r27 ($i64) $r27 $r1            // nothing live yet -- caller-save loop is empty
    + $r2 ($i64) $r0 3                 // arg0 = 3
    + $r3 ($i64) $r0 4                 // arg1 = 4
;
||
    $call add
;
||
    + $r1 ($i64) $r0 $r1              // capture RET (dest happened to land on $r1 too)
    + $r2 ($i64) $r26 -8               // &r, for the store below
;
||
    $st ($i32) [$r2 + 0] $r1          // r = (the call's result)
    + $r3 ($i64) $r26 -8
;
||
    $ld ($i32) $r4 [$r3 + 0]
;
||
    + $r1 ($i64) $r0 $r4
    ? ($i64) $r0 == $goto main_epilogue
;
main_epilogue:
||
    + $r27 ($i64) $r0 $r26
;
||
    $ld ($i64) $r26 [$r27 + 0]
    $return
;
\end{lstlisting}

Note the line \texttt{+ \$r1 = \$r0 + \$r1} in \texttt{main}: this looks like a no-op, and
on these exact values it is one --- but it is not special-cased anywhere. \texttt{main}
had nothing live at the call site, so the caller-save loop did nothing, and the register
allocator's ordinary linear-scan happened to hand out \texttt{\$r1} (the first register in
its free list) for the call's destination temp, which collides with \texttt{RET}
($=$\texttt{\$r1}) itself. The algorithm is correct either way --- it would be a real copy
into a \emph{different} register if anything else had been live first.

\subsection*{Caller-save actually firing: a second example}
\texttt{r = x + add(3, 4);} with \texttt{x} already assigned beforehand produces, at the
call site:
\begin{lstlisting}[caption={Caller-save in action (excerpt)}]
    $st ($i64) [$r26 + -96] $r3       // save x's value before clobbering $r3 for an argument
    + $r2 ($i64) $r0 3
;
||
    + $r3 ($i64) $r0 4                // $r3 reused for arg1 -- safe, x is already on the stack
    $call add
;
||
    + $r4 ($i64) $r0 $r1              // capture return value into a fresh register
    $ld ($i64) $r3 [$r26 + -96]       // restore x
;
||
    + $r5 ($i64) $r3 $r4              // r5 = x + add(3,4)
\end{lstlisting}

% ============================================================================
\section{Known Limits and Historical Failure Modes}
\label{sec:limits}
% ============================================================================

Four real, previously-confirmed defects sit directly in this machinery and are worth
recording alongside the algorithm itself, since each one demonstrates a way this exact
code path can go wrong if changed carelessly:
\begin{itemize}
  \item \textbf{Narrow-parameter bug (compiler-side, fixed).} The prologue's parameter
        store used to be a hardcoded \texttt{\$st (\$i64)} regardless of the parameter's
        real width, while reads correctly used the narrow width --- so any
        \texttt{int}/\texttt{short}/\texttt{char} parameter silently read back as garbage.
        Fixed by threading the real element width through \texttt{IRFuncBegin.params}
        (Section~\ref{sec:callee}'s width-correct store is the fix).
  \item \textbf{4-argument ceiling (compiler-side, fixed).} A 5th argument or parameter
        used to be silently dropped with no error. Now both \texttt{\_gen\_IRCall} and
        \texttt{\_gen\_IRFuncBegin} fail loudly at compile time instead.
  \item \textbf{\texttt{\$call} disassembler sign-extend bug (simulator-side, fixed).}
        Every \emph{backward} call (a callee defined above its caller in source order ---
        the ordinary C pattern) resolved to the wrong target entirely, because the
        disassembler sign-extended the wrong bit of the 25-bit jump field. Fixed in the
        simulator; confirmed present in the current authoritative build.
  \item \textbf{Instruction-memory size bug (simulator-side, fixed).} Functions whose
        compiled body exceeded 2048 words were silently truncated, which surfaced as
        \emph{this} machinery's symptoms (wrong return values) despite the bug having
        nothing to do with calling convention logic at all.
\end{itemize}
Indirect calls through a function pointer (\texttt{\$call \$rN}) are valid hardware
behaviour and fully implemented on the compiler side (\texttt{IRFuncAddr}/
\texttt{IRIndirectCall}), but are blocked end-to-end because no instruction exists to load
a function's absolute address into a register in the first place --- a hard assembler
limitation, not a bug in this calling convention.

% ============================================================================
\section{Summary}
% ============================================================================

A C function call is, by the time it reaches mcode, five mechanical steps: spill whatever
is live, load the fixed argument registers from the safety of the stack rather than from
registers that are about to be overwritten, call, capture \texttt{\$r1}, and reload
everything that was spilled. A function definition is two mechanical bookends around an
ordinary instruction stream: a prologue that links a new frame onto the old one and stores
incoming arguments at their real width, and a single epilogue, shared by every
\texttt{return} in the function, that unwinds the frame and hands control back. Nothing
here depends on what kind of function is being called or how many times it is called from
--- the same five-step/two-bookend pattern produced every line in
Section~\ref{sec:worked}'s worked example, and produces every function call in every
program this compiler has ever built.

\end{document}
